% 2013 重庆卷理科第 21 题：椭圆与圆 Q。
% 椭圆 x^2/16+y^2/8=1，图中按 3:4 缩放绘制（半长轴 3，半短轴、半焦距 3*sqrt2/2）。
% A、A' 为过左焦点 F1 的垂线与椭圆的交点；P、P' 为另一条垂线与椭圆的交点；
% 圆 Q 过 P、P' 且与椭圆在 P、P' 处相切，PQ⊥P'Q。
\begin{tikzpicture}[font=\small, line width=0.45pt, >=Stealth, scale=1]
  \def\aa{3}
  \def\bb{2.12132}
  \def\cc{2.12132}
  \def\yy{1.5}
  \def\px{2.44949}
  \def\py{1.22474}
  \def\qx{1.22474}
  \def\rr{1.73205}
  % 坐标轴
  \draw[->] (-3.75,0) -- (4.05,0) node[below] {$x$};
  \draw[->] (0,-2.65) -- (0,2.75) node[right] {$y$};
  \node[below left] at (0,0) {$O$};
  % 椭圆
  \draw (0,0) ellipse ({\aa} and {\bb});
  % 过左焦点 F1 的垂线
  \draw (-\cc,-\yy) -- (-\cc,\yy);
  \fill (-\cc,0) circle (0.6pt);
  \node[below left] at (-\cc,0) {$F_1$};
  \node[above left] at (-\cc,\yy) {$A$};
  \node[below left] at (-\cc,-\yy) {$A'$};
  % 过 P、P' 的垂线
  \draw (\px,-\py) -- (\px,\py);
  \node[above right] at (\px,\py) {$P$};
  \node[below right] at (\px,-\py) {$P'$};
  % 圆 Q
  \draw (\qx,0) circle ({\rr});
  \fill (\qx,0) circle (0.7pt);
  \node[below] at (\qx,0) {$Q$};
  \draw (\qx,0) -- (\px,\py);
  \draw (\qx,0) -- (\px,-\py);
\end{tikzpicture}
